\section{FPGA 原型与交付材料说明}

\subsection{FPGA 原型路线}

虽然赛方已经明确初赛阶段可以以文档和仿真为主，但本项目仍保留了完整的 FPGA pre-board 路线，
以证明方案具备原型实现可行性。当前原型目标平台为 Nexys A7-100T，主构建入口为
\texttt{impl50}，即 50MHz 冻结路径。

\subsection{当前 pre-board 已完成项}

截至当前版本，以下内容已经收口：

\begin{itemize}
    \item Vivado \texttt{impl50} 构建脚本、bitstream 路径和报告目录已冻结；
    \item 串口口径固定为 \texttt{115200 8N1}；
    \item 默认 payload 已与主工程脚本口径对齐；
    \item \texttt{IMEM\_OUTPUT\_REG=0} / \texttt{DMEM\_OUTPUT\_REG=0} 的近 FPGA 路径
    已通过 FPGA-like CoreMark probe 验证。
\end{itemize}

\begin{table}[H]
    \centering
    \caption{FPGA pre-board 当前可引用结论}
    \begin{tabular}{p{0.28\textwidth} p{0.52\textwidth}}
        \toprule
        项目 & 当前状态 \\
        \midrule
        bitstream 路径 & \texttt{project/YH\_rv\_cpu\_nexys\_a7\_100\_20p000.bit} \\
        报告目录 & \texttt{project/reports/clk\_20p000ns/} \\
        impl50 利用率 & \texttt{2556 LUT / 2170 FF / 4 BRAM / 0 DSP} \\
        impl50 时序 & \texttt{WNS = +5.599ns, WHS = +0.025ns} \\
        FPGA-like probe & \texttt{156442 cycles, 7.728811 CoreMark/MHz} \\
        \bottomrule
    \end{tabular}
\end{table}

\subsection{当前未完成项及其性质}

需要明确的是，当前未完成内容属于外部硬件条件阻塞，而不是设计说明书本身的缺口：

\begin{enumerate}
    \item 实体板卡尚未到位，无法完成 UART/LED 实板闭环；
    \item 当前 XDC 仍是 pre-board 约束口径，缺最终板级 I/O delay signoff；
    \item 因此项目可以诚实地宣称 ``pre-board closure 已完成''，但不能宣称
    ``real-board signoff 已完成''。
\end{enumerate}

\FigurePlaceholder{FPGA pre-board 流程与证据路径图占位}

\subsection{为什么当前阶段不优先推进上板}

结合赛方答疑和项目主线状态，可以得出明确结论：初赛阶段的硬交付重点是设计/技术文档，
而不是强行把所有外部依赖都压缩到当前窗口内完成。因此，本项目采取的策略是：

\begin{itemize}
    \item 先把源码、文档、性能与验证口径收敛成可提交、可答辩的工程基线；
    \item 将 FPGA pre-board 结果作为增强支撑材料保留；
    \item 等实体板卡到位后，再执行 UART、LED、boot banner 和板级约束闭环。
\end{itemize}

\subsection{初赛提交材料组织方式}

当前比赛相关材料已经统一迁入 \texttt{01-项目管理/05-汇报与提交材料/}，
避免出现源码主线、汇报材料和赛题要求三套口径分别维护的问题。初赛目录中的主要文件分工如下：

\begin{table}[H]
    \centering
    \caption{初赛材料目录分工}
    \begin{tabular}{p{0.30\textwidth} p{0.48\textwidth}}
        \toprule
        文件/目录 & 作用 \\
        \midrule
        \texttt{README.md} & 初赛材料入口与统一说明 \\
        \texttt{初赛提交清单.md} & 标识哪些材料 ready、哪些仍 blocked \\
        \texttt{冻结基线与证据索引.md} & 对外引用结果与证据路径索引 \\
        \texttt{答辩汇报口径.md} & 汇报、PPT、答辩可直接复用的话术 \\
        \texttt{latex/} & 本文档的 LaTeX 源码与后续插图接口 \\
        \bottomrule
    \end{tabular}
\end{table}

\subsection{当前交付建议}

对于初赛阶段，建议将本说明书作为主交付文档之一，同时配合以下材料共同组成提交包：

\begin{itemize}
    \item \texttt{YH\_rv\_cpu/doc/技术文档.md}
    \item \texttt{YH\_rv\_cpu/README.md}
    \item \texttt{YH\_rv\_cpu/doc/YH\_rv\_cpu\_handoff.md}
    \item \texttt{01-项目管理/01-赛题要求/七星微赛题答疑整理.md}
    \item 当前目录中的提交清单、证据索引和答辩口径文件
\end{itemize}

\clearpage
